% GENERATED FILE. Edit canonical lessons, then run npm run generate:books.
% canonical-source-hash: fnv1a64:063de13f8c535ea8
% canonical-lessons: ZH-C284-zhantai, ZH-C284-ticket, ZH-W284-zuo, ZH-C284-dancheng, ZH-C284-zuowei, ZH-C284-tuipiao, ZH-C284-v-zai, ZH-C284-jia-measure, ZH-C284-practice

\chapter{Two Tickets}
\label{ch:zh-two-tickets}
\hlchaptermodality{284}

\begin{chapteropening}
\textbf{By the end of this chapter:} \emph{I can buy a ticket and say how I am travelling.}

Saying how you travel with \zh{坐} … \zh{去}, buying and returning tickets, platforms, one-way tickets and seats, \zh{在} after a verb for where something lands, and \zh{家} for counting shops, built on one new written character: \zh{坐}.
\end{chapteropening}

\section[zhàntái]{zhàntái — a platform}

\label{lesson:ZH-C284-zhantai}



\begin{quote}
A new chapter, and it opens at the station, where the trains stand.

\begin{itemize}
\raggedright
  \item \emph{Recall:} say \zh{去书店怎么走} — \emph{qù shūdiàn zěnme zǒu} — \textbf{R1}, the chapter you have finished
  \item \emph{Recall:} say \zh{路过} — \emph{lùguò}, to pass by — \textbf{R2}
  \item \emph{Recall:} say \zh{火车} — \emph{huǒchē}, a train — \textbf{R3}, durable
  \item \emph{Recall:} say \zh{火车站} — \emph{huǒchēzhàn}, a station — \textbf{R3}
\end{itemize}
\end{quote}

\subsection*{You'll want to know: zhàntái}
\emph{zhàntái} — \textquotedblleft{}a platform\textquotedblright{}, at a station.

\begin{quote}
\emph{huǒchē zài èr hào zhàntái} — The train is at platform two.
\end{quote}

Its second character comes later; for now it is a word for the ear.

\begin{cousinweb}
\emph{zhàntái} is built from \zh{站} and tái 'a raised platform'. The first is inherited from Middle Chinese \emph{treamH} (Guangyun, in Baxter's transcription), which the dictionary gives as the everyday word for standing alone. tái is written with a character whose traditional form is \emph{doj}, glossed as a raised square of earth; the simplified character is also an old one in its own right, with the readings \emph{yi} and \emph{thoj} and other senses.
\end{cousinweb}

\begin{grammarlens}[title={What you've built}]
\emph{zhàntái} holds the \zh{站} of \zh{火车站}: the station, and the platform you stand on in it.

A number comes before \zh{号}: \emph{èr hào zhàntái}, platform two.
\end{grammarlens}

\subsection*{Your turn}
\begin{itemize}
\raggedright
  \item \emph{Say it:} \emph{zhàntái}
  \item \emph{Say it:} \emph{zhàntái}, once more
  \item \emph{Say it:} \emph{huǒchē zài èr hào zhàntái} — The train is at platform two.
  \item \emph{Recall:} say \zh{火车站}, then \emph{zhàntái}
\end{itemize}

\begin{tcolorbox}[breakable,colback=teal!4,colframe=teal!35!black,title={Before you move on}]
What does \emph{zhàntái} mean? (\textquotedblleft{}A platform\textquotedblright{}.) How do you say platform two? (\emph{èr hào zhàntái}.)
\end{tcolorbox}
\section[\texorpdfstring{\zh{坐火车去} (zuò huǒchē qù)}{坐火车去 (zuò huǒchē qù)}]{\zh{坐火车去} — going by train}

\label{lesson:ZH-C284-ticket}



\begin{quote}
A verb you have said since the early chapters, and the ticket that goes with it.

\begin{itemize}
\raggedright
  \item \emph{Recall:} say \emph{zhàntái}, a platform — \textbf{R1}, one lesson back
  \item \emph{Recall:} say \emph{zuò}, to sit; to travel by — \textbf{R3}, durable
  \item \emph{Recall:} say \zh{票} — \emph{piào}, a ticket — \textbf{R3}
  \item \emph{Recall:} say \emph{zhāng}, the measure word for flat things — \textbf{R3}
\end{itemize}
\end{quote}

\begin{grammarlens}[title={Grammar Lens: \zh{坐} and a vehicle, then \zh{去}}]
\begin{quote}
\textbf{\zh{坐} and a vehicle, then \zh{去} or \zh{来}: that is how you travel.}
\end{quote}

\begin{quote}
\textbf{\zh{我们坐火车去}} \emph{wǒmen zuò huǒchē qù} — We're going by train. \textbf{\zh{我们坐八点的火车}} \emph{wǒmen zuò bā diǎn de huǒchē} — We're taking the eight o'clock train.
\end{quote}

\zh{坐} is the \emph{zuò} you know, to sit and to travel by; the next lesson puts it in your hand. The vehicle comes between \zh{坐} and the verb of going.

Tickets are counted with \emph{zhāng}, as flat things are: \emph{mǎi liǎng zhāng qù Běijīng de piào}, two tickets to Beijing.
\end{grammarlens}

\subsection*{Your turn}
Say these aloud:

\begin{itemize}
\raggedright
  \item \zh{我们坐火车去} — \emph{wǒmen zuò huǒchē qù} — We're going by train.
  \item \zh{我们坐八点的火车} — \emph{wǒmen zuò bā diǎn de huǒchē} — We're taking the eight o'clock train.
  \item \emph{mǎi liǎng zhāng qù Běijīng de piào} — Two tickets to Beijing, please.
  \item how you get to work or school, with \zh{坐} … \zh{去}
\end{itemize}

\begin{tcolorbox}[breakable,colback=teal!4,colframe=teal!35!black,title={Before you move on}]
Where does the vehicle go? (\textbf{Between \zh{坐} and \zh{去}}.) What counts tickets? (\emph{zhāng}, as for flat things.)
\end{tcolorbox}
\section[\texorpdfstring{\zh{坐} (zuò)}{坐 (zuò)}]{\zh{坐} — \zh{人} twice, then \zh{土}}

\label{lesson:ZH-W284-zuo}



\begin{quote}
The sentence first, then the character it was waiting for.

\begin{itemize}
\raggedright
  \item \emph{Recall:} say \zh{我们坐火车去} — \emph{wǒmen zuò huǒchē qù} — \textbf{R1}, one lesson back, before it decays
\end{itemize}

\begin{quote}
\zh{人} \zh{从} \zh{在}
\end{quote}

\emph{Read it:} each one aloud, with nothing said first — \emph{rén}, \emph{cóng}, \emph{zài}
\end{quote}

\subsection*{Script — \zh{人} twice, then \zh{土}}
\hlblockfigure{\detokenize{figures/ZH-W284-zuo-filmstrip.pdf}}{How \zh{坐} is written, stroke by stroke}

\begin{quote}
\zh{坐}
\end{quote}

\textbf{Three pieces, and you write all of them already.}

\textbf{\zh{人}, twice, side by side}, the left one first. Each is a left-falling stroke and then a dot: neither has the long right-falling stroke of \zh{人} standing alone.

\textbf{\zh{土}, three strokes}, under them, as at the lower right of \zh{在}: the short horizontal; then the vertical, from the top between the two \zh{人}, down through it; then the long horizontal along the foot.

2 + 2 + 3 = seven.

\subsection*{Your turn}
\begin{itemize}
\raggedright
  \item \emph{Write it:} \zh{坐}, all seven, with the model in front of you
  \item \emph{Cover:} write it again from memory, and count the pieces, not the strokes
\end{itemize}

\begin{tcolorbox}[breakable,colback=teal!4,colframe=teal!35!black,title={Before you move on}]
How many strokes in \zh{坐}? (\textbf{Seven}.) What shrinks to a dot in each \zh{人}? (\textbf{The second stroke}.) Where does the vertical of \zh{土} start? (\textbf{At the top, between the two \zh{人}}.)

Source: \href{https://raw.githubusercontent.com/chanind/hanzi-writer-data/68d10a4b21150cae5e1ebbd223eed289cf32d90c/data/\%E5\%9D\%90.json}{Hanzi Writer Data \zh{坐}, pinned stroke paths and medians}.
\end{tcolorbox}
\section[dānchéng]{dānchéng — one way}

\label{lesson:ZH-C284-dancheng}



\begin{quote}
Before the new one: the train you are taking, and how you count its tickets.

\begin{itemize}
\raggedright
  \item \emph{Read it:} \zh{我们坐八点的火车}, aloud — \emph{wǒmen zuò bā diǎn de huǒchē}
  \item \emph{Recall:} say \emph{mǎi liǎng zhāng piào}, two tickets — \textbf{R2}
\end{itemize}
\end{quote}

\subsection*{You'll want to know: dānchéng}
\emph{dānchéng} — \textquotedblleft{}one way\textquotedblright{}, of a ticket or a trip.

\begin{quote}
\emph{yì zhāng dānchéng piào} — a one-way ticket
\end{quote}

\begin{quote}
\emph{dānchéng háishi wǎngfǎn} — One way or return?
\end{quote}

\emph{wǎngfǎn}, there and back, is the other choice. Its characters come later; for now \emph{dānchéng} is a word for the ear.

\begin{cousinweb}
\emph{dānchéng} is built from dān 'single' and chéng, both inherited from Middle Chinese, \emph{tan} and \emph{drjeng} (Guangyun, in Baxter's transcription). The dictionary glosses the first 'single, not double' and the second 'a term; a limit'; 'one way' for the pair, a single stage of a journey, is the word's own sense.
\end{cousinweb}

\begin{grammarlens}[title={What you've built}]
At the window you are asked \emph{dānchéng háishi wǎngfǎn?}, one way or return. \zh{还是} is the or of a question.

Answer with the word alone: \emph{dānchéng}.
\end{grammarlens}

\subsection*{Your turn}
\begin{itemize}
\raggedright
  \item \emph{Say it:} \emph{dānchéng}
  \item \emph{Say it:} \emph{dānchéng}, once more
  \item \emph{Say it:} \emph{yì zhāng dānchéng piào} — a one-way ticket
  \item \emph{Say it:} \emph{dānchéng háishi wǎngfǎn} — One way or return?
  \item \emph{Recall:} answer the question: \emph{dānchéng}
\end{itemize}

\begin{tcolorbox}[breakable,colback=teal!4,colframe=teal!35!black,title={Before you move on}]
What does \emph{dānchéng} mean? (\textquotedblleft{}One way\textquotedblright{}.) What is the other choice? (\emph{wǎngfǎn}, there and back.)
\end{tcolorbox}
\section[zuòwèi]{zuòwèi — a seat}

\label{lesson:ZH-C284-zuowei}



\begin{quote}
Before the new one: last lesson's word, and the character you wrote this chapter.

\begin{itemize}
\raggedright
  \item \emph{Recall:} say \emph{dānchéng}, one way — \textbf{R1}, one lesson back
\end{itemize}

\begin{quote}
\zh{坐}
\end{quote}

\emph{Read it:} it aloud, with nothing said first — \emph{zuò}
\end{quote}

\subsection*{You'll want to know: zuòwèi}
\emph{zuòwèi} — \textquotedblleft{}a seat\textquotedblright{}.

\begin{quote}
\emph{zhèr yǒu zuòwèi ma} — Is there a seat here?
\end{quote}

Its characters come later. The first is not the \zh{坐} you write but one that shares its sound and holds it under a roof.

\begin{cousinweb}
\emph{zuòwèi} is built from zuò 'a seat' and wèi 'a place', both inherited from Middle Chinese, \emph{dzwaH} and \emph{hwijH} (Guangyun, in Baxter's transcription), glossed 'the stand of a bed' and 'a position; a rank'.
\end{cousinweb}

\begin{grammarlens}[title={What you've built}]
\zh{坐} is the verb, to sit; \emph{zuòwèi} is the noun, the seat itself: \emph{wǒ de zuòwèi zài nǎr?}, where's my seat?
\end{grammarlens}

\subsection*{Your turn}
\begin{itemize}
\raggedright
  \item \emph{Say it:} \emph{zuòwèi}
  \item \emph{Say it:} \emph{zuòwèi}, once more
  \item \emph{Say it:} \emph{zhèr yǒu zuòwèi ma} — Is there a seat here?
  \item \emph{Say it:} \emph{wǒ de zuòwèi zài nǎr} — Where's my seat?
  \item \emph{Recall:} say \zh{坐}, then \emph{zuòwèi}: to sit, then the seat
\end{itemize}

\begin{tcolorbox}[breakable,colback=teal!4,colframe=teal!35!black,title={Before you move on}]
What does \emph{zuòwèi} mean? (\textquotedblleft{}A seat\textquotedblright{}.) Is its first character the \zh{坐} you write? (\textbf{No: the same sound, another character}.)
\end{tcolorbox}
\section[tuìpiào]{tuìpiào — to return a ticket}

\label{lesson:ZH-C284-tuipiao}



\begin{quote}
Before the new one: last lesson's word, and one you use at a hotel desk.

\begin{itemize}
\raggedright
  \item \emph{Recall:} say \emph{zuòwèi}, a seat — \textbf{R1}, one lesson back
  \item \emph{Recall:} say \emph{tuìfáng}, to check out — \textbf{R3}, durable
\end{itemize}

\begin{quote}
\zh{票}
\end{quote}

\emph{Read it:} it aloud, with nothing said first — \emph{piào}
\end{quote}

\subsection*{You'll want to know: tuìpiào}
\emph{tuìpiào} — \textquotedblleft{}to return a ticket\textquotedblright{}, for your money back.

\begin{quote}
\emph{wǒ xiǎng tuìpiào} — I'd like to return my ticket.
\end{quote}

Its first character comes later; for now it is a word for the ear.

\begin{cousinweb}
\emph{tuìpiào} is built from tuì 'to give back' and \zh{票} 'a ticket'. The first is inherited from Middle Chinese \emph{thwojH} (Guangyun, in Baxter's transcription), glossed 'to draw back'; the dictionary gives no form for the second in its ticket sense, so that part's history is not settled.
\end{cousinweb}

\begin{grammarlens}[title={What you've built}]
You know \emph{tuìfáng}, to give back a hotel room; \emph{tuìpiào} gives back a ticket the same way.

It is a verb on its own, with the ticket inside it: \emph{wǒ xiǎng tuìpiào}.
\end{grammarlens}

\subsection*{Your turn}
\begin{itemize}
\raggedright
  \item \emph{Say it:} \emph{tuìpiào}
  \item \emph{Say it:} \emph{tuìpiào}, once more
  \item \emph{Say it:} \emph{wǒ xiǎng tuìpiào} — I'd like to return my ticket.
  \item \emph{Recall:} say \emph{tuìfáng}, then \emph{tuìpiào}
\end{itemize}

\begin{tcolorbox}[breakable,colback=teal!4,colframe=teal!35!black,title={Before you move on}]
What does \emph{tuìpiào} mean? (\textquotedblleft{}To return a ticket\textquotedblright{}.) Which word gives back a hotel room? (\emph{tuìfáng}.)
\end{tcolorbox}
\section[\texorpdfstring{\zh{坐在这儿} (zuò zài zhèr)}{坐在这儿 (zuò zài zhèr)}]{\zh{坐在这儿} — \zh{在} after the verb}

\label{lesson:ZH-C284-v-zai}



\begin{quote}
A word you have used since the first chapters, in a new place in the sentence.

\begin{itemize}
\raggedright
  \item \emph{Recall:} say \emph{tuìpiào}, to return a ticket — \textbf{R1}, one lesson back
  \item \emph{Recall:} say \zh{我在家} — \emph{wǒ zài jiā}, I'm at home — \textbf{R3}, durable
  \item \emph{Recall:} say \emph{zhù}, to live — \textbf{R3}
  \item \emph{Recall:} say \emph{fàng}, to put — \textbf{R3}
\end{itemize}
\end{quote}

\begin{grammarlens}[title={Grammar Lens: \zh{在} after the verb, for where the action lands}]
\begin{quote}
\textbf{\zh{在} after a verb says where the action lands, or where it stays.}
\end{quote}

\begin{quote}
\textbf{\zh{你坐在这儿}} \emph{nǐ zuò zài zhèr} — You sit here. \emph{tā zhù zài xuéxiào lǐ} — He lives at the school.
\end{quote}

Keep it apart from \zh{在} before the verb: \zh{我在家吃饭} says where the eating happens. After the verb, \zh{在} marks where you end up sitting, living or putting something: \emph{fàng zài zhèr}, put it here.
\end{grammarlens}

\subsection*{Your turn}
Say these aloud:

\begin{itemize}
\raggedright
  \item \zh{你坐在这儿} — \emph{nǐ zuò zài zhèr} — You sit here.
  \item \emph{tā zhù zài xuéxiào lǐ} — He lives at the school.
  \item \emph{fàng zài zhèr} — Put it here.
  \item where you live, with \emph{zhù zài}
\end{itemize}

\begin{tcolorbox}[breakable,colback=teal!4,colframe=teal!35!black,title={Before you move on}]
Where does \zh{在} go in \zh{你坐在这儿}? (\textbf{After the verb}.) What does it say there? (\textbf{Where the sitting lands}.)
\end{tcolorbox}
\section[\texorpdfstring{\zh{一家书店} (yì jiā shūdiàn)}{一家书店 (yì jiā shūdiàn)}]{\zh{一家书店} — counting shops with \zh{家}}

\label{lesson:ZH-C284-jia-measure}



\begin{quote}
A word you wrote in the first chapters, doing a new job.

\begin{itemize}
\raggedright
  \item \emph{Read it:} \zh{你坐在这儿}, aloud — \emph{nǐ zuò zài zhèr}
  \item \emph{Recall:} say \zh{书店} — \emph{shūdiàn}, a bookshop — \textbf{R3}, durable
  \item \emph{Recall:} say \emph{bīnguǎn}, a hotel — \textbf{R3}
\end{itemize}
\end{quote}

\begin{grammarlens}[title={Grammar Lens: \zh{家} for counting shops}]
\begin{quote}
\textbf{\zh{家} after a number, or after \zh{这} or \zh{那}, counts shops, hotels and restaurants.}
\end{quote}

\begin{quote}
\textbf{\zh{这家书店很好}} \emph{zhè jiā shūdiàn hěn hǎo} — This bookshop is very good. \textbf{\zh{那家饭店的菜很好吃}} \emph{nà jiā fàndiàn de cài hěn hǎochī} — The food at that restaurant is delicious.
\end{quote}

\zh{一家书店}, a bookshop; \emph{liǎng jiā bīnguǎn}, two hotels. It is the \zh{家} of home, now counting places that do business. \zh{个} does not count them.
\end{grammarlens}

\subsection*{Your turn}
Say these aloud:

\begin{itemize}
\raggedright
  \item \zh{这家书店很好} — \emph{zhè jiā shūdiàn hěn hǎo} — This bookshop is very good.
  \item \zh{那家饭店的菜很好吃} — \emph{nà jiā fàndiàn de cài hěn hǎochī} — The food at that restaurant is delicious.
  \item \emph{liǎng jiā bīnguǎn} — two hotels
  \item the name of a shop you like, with \zh{这家}
\end{itemize}

\begin{tcolorbox}[breakable,colback=teal!4,colframe=teal!35!black,title={Before you move on}]
What does \zh{家} count? (\textbf{Shops, hotels and restaurants}.) Does \zh{个} count them? (\textbf{No}.)
\end{tcolorbox}
\section[Practice]{\zh{我们坐八点的火车}}

\label{lesson:ZH-C284-practice}



\begin{quote}
The chapter closes here, at the ticket window.

\begin{itemize}
\raggedright
  \item \emph{Read it:} \zh{这家书店很好}, aloud — \emph{zhè jiā shūdiàn hěn hǎo}
  \item \emph{Recall:} say \emph{tuìpiào}, to return a ticket — \textbf{R2}
  \item \emph{Recall:} say \emph{zuòwèi}, a seat — \textbf{R2}
\end{itemize}
\end{quote}

\subsection*{The exchange}
\begin{quote}
\textbf{A:} \emph{wǒ yào liǎng zhāng qù Běijīng de piào} — Two tickets to Beijing, please.
\end{quote}

\begin{quote}
\textbf{B:} \emph{dānchéng háishi wǎngfǎn} — One way or return?
\end{quote}

\begin{quote}
\textbf{A:} \emph{dānchéng, wǒmen zuò bā diǎn de huǒchē} — One way. We're taking the eight o'clock train.
\end{quote}

\begin{quote}
\textbf{B:} \emph{bā diǎn de huǒchē zài èr hào zhàntái} — The eight o'clock is at platform two.
\end{quote}

\begin{quote}
\textbf{A:} \emph{yǒu zuòwèi ma} — Are there seats?
\end{quote}

\begin{quote}
\textbf{B:} \zh{有}
\end{quote}

\begin{quote}
\emph{yǒu} — There are.
\end{quote}

\begin{quote}
\textbf{A:} \zh{车站里有一家书店吗}
\end{quote}

\begin{quote}
\emph{chēzhàn lǐ yǒu yì jiā shūdiàn ma} — Is there a bookshop in the station?
\end{quote}

\begin{quote}
\textbf{B:} \emph{yǒu, zài zhàntái qiánbian} — There is, in front of the platform.
\end{quote}

The answers are short: \emph{dānchéng}, \zh{有}. The station's bookshop is counted with \zh{家}, and the platform has a number before \zh{号}.

\begin{tcolorbox}[breakable,skin=enhanced,colback=orange!6,colframe=orange!45!black,boxrule=0.5pt,arc=1mm,left=6pt,right=6pt,top=4pt,bottom=4pt,fonttitle=\bfseries,title={Read this}]
\emph{Read it:} the exchange through once without stopping, then answer the question below
\end{tcolorbox}

\begin{grammarlens}[title={What you've built — getting on a train}]
Everything you need between the window and your seat:

\begin{itemize}
\raggedright
  \item \textbf{\zh{我们坐火车去}} \emph{wǒmen zuò huǒchē qù} — how you travel.
  \item \emph{mǎi liǎng zhāng piào}, \emph{dānchéng}, \emph{tuìpiào} — buying, and giving back.
  \item \emph{èr hào zhàntái}, \emph{zuòwèi} — where to go, and where to sit.
  \item \textbf{\zh{你坐在这儿}} \emph{nǐ zuò zài zhèr} — where you land.
  \item \textbf{\zh{一家书店}} \emph{yì jiā shūdiàn} — a shop, counted.
\end{itemize}

\textbf{Listen.} \emph{Nǐmen zuò jǐ diǎn de huǒchē?}

\textbf{Speak.} Buy a ticket for a trip you would like to take: where, how many, one way or return.

\emph{Write it:} \zh{坐火车} — \zh{人} twice, then the foot of \zh{在}
\end{grammarlens}

\subsection*{Your turn}
\begin{itemize}
\raggedright
  \item \emph{Say it:} \zh{我们坐八点的火车} — \emph{wǒmen zuò bā diǎn de huǒchē} — We're taking the eight o'clock train.
  \item \emph{Say it:} \emph{wǒ yào liǎng zhāng qù Běijīng de piào} — Two tickets to Beijing, please.
  \item \emph{Say it:} the whole exchange, both parts
  \item \emph{Write it:} \zh{我们坐火车去}
\end{itemize}

\begin{tcolorbox}[breakable,colback=teal!4,colframe=teal!35!black,title={Before you move on}]
Say you're going by train. (\textbf{\zh{我们坐火车去}}.) Ask for one-way. (\emph{dānchéng}.) Tell someone to sit here. (\textbf{\zh{你坐在这儿}}.)
\end{tcolorbox}
